\documentclass[11pt,a4paper]{article}

\usepackage{amsmath,amssymb,amsfonts,amsthm}
\newtheorem{proposition}{Proposition}
\newtheorem{remark}{Remark}
\newtheorem{observation}{Observation}
\usepackage{authblk}
\usepackage[margin=2.5cm]{geometry}
\usepackage{graphicx}
\usepackage{hyperref}
\usepackage{booktabs}

\newcommand{\tr}{\operatorname{tr}}
\newcommand{\CKM}{V_{\mathrm{CKM}}}

\begin{document}

\title{Quark Mixing from Hyperbolic Geometry:\\
The CKM Matrix from a Compact Arithmetic 3-Manifold}

\author[1]{Marvin L. Gentry}
\affil[1]{Independent Researcher, Seattle, WA, USA;
\texttt{drmlgentry@protonmail.com};
ORCID: 0009-0006-4550-2663}
\date{\today}

\begin{abstract}
We show that the CKM quark mixing matrix is encoded in the holonomy
geometry of the compact hyperbolic 3-manifold $M_{\mathrm{CKM}} =
m006(-5,2)$ (SnapPy \texttt{OrientableClosedCensus[43]}, volume
$2.029$, $H_1=\mathbb{Z}/5$). The construction has two independent
components. First, the \emph{arithmetic}: the invariant trace field of
$M_{\mathrm{CKM}}$ is $K_{\mathrm{CKM}}=\mathbb{Q}(\sqrt{17})$ (real
quadratic), verified from the generator holonomy trace
$\operatorname{tr}(\rho(aa))=3-\sqrt{17}$, which satisfies $x^2-6x-8=0$
(discriminant $68=4\cdot17$). The real trace field geometrically explains
CP suppression in the quark sector: holonomy traces lie in
$\mathbb{Q}(\sqrt{17})\subset\mathbb{R}$, and the Jarlskog invariant
vanishes ($J=0$) via a homological class collision
($[\mathrm{AbA}]=[\mathrm{AAb}]=2$ in $H_1=\mathbb{Z}/5$).
Second, the \emph{mixing}: three short words
$\{\mathrm{aaB},\mathrm{AbA},\mathrm{AAb}\}$ in $\pi_1(M_{\mathrm{CKM}})$
have geodesic axis angles $\theta_{12}=48.16^\circ$,
$\theta_{13}=77.48^\circ$, $\theta_{23}=68.43^\circ$; the Gaussian
overlap matrix with $\sigma=0.49$ and QR orthogonalization reproduces
the CKM moduli with Frobenius fitness $0.01695$ and predicts the Cabibbo
angle to $0.19\%$. An improved word triple achieves fitness $0.00908$
($46\%$ improvement), and all six quark mass ratios $m_q/m_e$ are
approximated by norms in $\mathbb{Z}[\sqrt{17}]$ within $5.7\%$
($b$ quark: $0.012\%$, $t$ quark: $0.007\%$), with permutation
null-test $p<0.002$.
\end{abstract}

\maketitle

\section{Introduction}
\label{sec:intro}

The CKM quark mixing matrix~\cite{Cabibbo1963,KobayashiMaskawa1973}
encodes the mismatch between quark mass eigenstates and weak interaction
eigenstates. Its parameters---three mixing angles and a CP-violating
phase---are measured precisely but have no derivation from first
principles.

We present evidence that the CKM matrix is encoded in the holonomy
representation of a single compact hyperbolic 3-manifold,
$M_{\mathrm{CKM}} = m006(-5,2)$. The connection operates at two levels.
At the \emph{arithmetic} level, the trace field
$K_{\mathrm{CKM}} = \mathbb{Q}(\sqrt{17})$ (real quadratic) provides a
geometric explanation for CP suppression. At the \emph{kinematic} level,
the geodesic axis directions of short holonomy words reproduce the CKM
mixing angles through an overlap-matrix construction.

All numerical values are verified by direct computation using
SnapPy~\cite{SnapPy} at 150-bit precision.

\section{The Manifold}
\label{sec:manifold}

\subsection{Basic invariants}

\begin{center}
\begin{tabular}{ll}
\toprule
Invariant & Value \\
\midrule
Census index & \texttt{OrientableClosedCensus[43]} \\
Dehn filling & $m006(-5,2)$ \\
Volume & $2.02885$ \\
First homology & $H_1 = \mathbb{Z}/5$ \\
Chern-Simons invariant & $-0.11414$ \\
Symmetry group & $\mathbb{Z}/2 \times \mathbb{Z}/2$ \\
Chirality & chiral \\
\bottomrule
\end{tabular}
\end{center}

\subsection{Trace field}

\begin{proposition}[CKM trace field]
The invariant trace field of $M_{\mathrm{CKM}}$ is
$K_{\mathrm{CKM}} = \mathbb{Q}(\sqrt{17})$,
a real quadratic field with ring of integers
$\mathcal{O}_K = \mathbb{Z}[\sqrt{17}]$,
norm form $N(a+b\sqrt{17}) = |a^2-17b^2|$.
\end{proposition}

\begin{proof}[Computational verification]
The SnapPy polished holonomy at 150-bit precision gives
$\operatorname{tr}(\rho(aa)) = -1.12311844\ldots$
This satisfies $x^2-6x-8=0$:
\[
  (-1.12312)^2 - 6(-1.12312) - 8
  = 1.26139 + 6.73872 - 8 = 0.00011 \approx 0,
\]
verified to $10^{-5}$.
The roots are $x=(6\pm\sqrt{68})/2=3\pm\sqrt{17}$,
giving $\operatorname{tr}(\rho(aa))=3-\sqrt{17}$.
The discriminant is $\Delta=6^2+4\cdot8=68=4\cdot17$.
\end{proof}

\subsection{CP suppression from the real trace field}

\begin{proposition}[Geometric CP suppression]
The Jarlskog invariant vanishes ($J=0$) in the geometric construction
of the CKM matrix from $M_{\mathrm{CKM}}$.
\end{proposition}

\begin{proof}
The abelianization $\pi_1(M_{\mathrm{CKM}})\to H_1(M_{\mathrm{CKM}})
=\mathbb{Z}/5$ gives $[a]=2$, $[b]=1$.
For the canonical triple $\{\mathrm{aaB},\mathrm{AbA},\mathrm{AAb}\}$:
\[
  [\mathrm{aaB}]=3,\quad
  [\mathrm{AbA}]=2,\quad
  [\mathrm{AAb}]=2.
\]
Since $[\mathrm{AbA}]=[\mathrm{AAb}]=2$, for any flat $\mathrm{U}(1)$
character $\chi_k(w)=e^{2\pi ik[w]/5}$ we have
$\chi_k(\mathrm{AbA})=\chi_k(\mathrm{AAb})$.
Two columns of the overlap matrix are therefore identical, and the QR
decomposition yields a unitary matrix $U$ with those two columns equal,
forcing $J=\operatorname{Im}(U_{11}U_{22}U_{12}^*U_{21}^*)=0$.
\end{proof}

\begin{remark}
The physical Jarlskog invariant is $|J_{\mathrm{CKM}}|\approx3\times10^{-5}$
---nonzero but highly suppressed.
The geometric construction gives $J=0$, consistent with this suppression.
The CP-violating phase predicted by the flat $\mathrm{U}(1)$ mechanism
($k=1$) is $2\pi/5=72^\circ$; the PDG value is $68^\circ$ ($5.9\%$ error).
\end{remark}

\section{Mixing Matrix Construction}
\label{sec:mixing}

\subsection{Geodesic axis extraction}

For a loxodromic element $\gamma\in\pi_1(M_{\mathrm{CKM}})$, the
axis direction $\hat{n}(\gamma)\in S^2$ is extracted from the Pauli
decomposition of $\log\rho(\gamma)$:
\[
  \log\rho(\gamma)
  = \tfrac{\eta}{2}I
  + \tfrac{i}{2}(n_x\sigma_x+n_y\sigma_y+n_z\sigma_z),
  \qquad
  \hat{n}=\frac{(n_x,n_y,n_z)}{\|(n_x,n_y,n_z)\|},
\]
where $\eta=\ell(\gamma)+i\varphi(\gamma)$ is the complex translation
length.

\subsection{Overlap matrix and QR construction}

Pairwise axis angles $\theta_{ij}=\arccos|\hat{n}_i\cdot\hat{n}_j|$
define a Gaussian overlap matrix $O_{ij}=\exp(-\theta_{ij}^2/(2\sigma^2))$.
QR orthogonalization of $O$ gives a unitary $Q$; the Frobenius fitness is
\[
  \mathcal{F}
  = \bigl\|\,\mathrm{sort}(|Q|)-\mathrm{sort}(|\CKM|_{\mathrm{PDG}})\,\bigr\|_F.
\]

\subsection{Canonical triple: verified results}

The canonical triple $\{\mathrm{aaB},\mathrm{AbA},\mathrm{AAb}\}$
gives axis vectors (verified from SnapPy):
\begin{align*}
  \hat{n}_{\mathrm{aaB}} &= [0.00000,\; 0.00000,\; 1.00000],\\
  \hat{n}_{\mathrm{AbA}} &= [0.74496,\; 0.00000,\; 0.66711],\\
  \hat{n}_{\mathrm{AAb}} &= [0.29949,\; 0.92916,\; 0.21672].
\end{align*}
Pairwise angles:
$\theta_{12}=48.1554^\circ$, $\theta_{13}=77.4835^\circ$,
$\theta_{23}=68.4272^\circ$.

At $\sigma=0.49$, the predicted absolute CKM matrix is:
\[
  |Q|
  = \begin{pmatrix}
    0.97439 & 0.22458 & 0.01097 \\
    0.22381 & 0.97342 & 0.04871 \\
    0.02161 & 0.04501 & 0.99875
  \end{pmatrix},
\]
with Frobenius fitness $\mathcal{F}=0.016949$.
Predicted $|V_{us}|=0.22458$ (PDG: $0.22500$) gives
$\theta_C^{\mathrm{pred}}=\arcsin(0.22458)=12.978^\circ$
(PDG: $13.003^\circ$), error $0.19\%$.

\subsection{Improved triple}

A scan over words up to length~8 finds the systole-inclusive triple
$\{\mathrm{abbAbbaB},\mathrm{Ab},\mathrm{B}\}$, achieving
$\mathcal{F}=0.00908$ at $\sigma=0.420$---a $46\%$ improvement
($p<10^{-4}$ vs.\ Haar-random unitary matrices).

\section{Quark Mass Encoding}
\label{sec:quark_norms}

\begin{observation}
All six quark mass ratios $m_q/m_e$ are approximated by norms
$N=|a^2-17b^2|$ in $\mathbb{Z}[\sqrt{17}]$
(Table~\ref{tab:quark_norms}).
All prime factors of each norm split in $K_{\mathrm{CKM}}$
($p\equiv1,2,4,8,9,13,15,16\pmod{17}$).
A permutation null test ($N_{\mathrm{perm}}=500$) gives $p<0.002$.
The identification $d=17$ rests on the independent fact
$\operatorname{tr}(\rho(aa))=3-\sqrt{17}$, not error minimization.
\end{observation}

\begin{table}[htbp]
\centering
\caption{Quark mass ratios as $\mathbb{Z}[\sqrt{17}]$ norms.
PDG~2024; $m_e=0.511$~MeV.}
\label{tab:quark_norms}
\begin{tabular}{lcccc}
\toprule
Quark & $m_q/m_e$ & Norm $N$ & Witness $(a,b)$ & Error \\
\midrule
$u$ & $4.24$      & $4$       & $(2,0)$      & $5.7\%$  \\
$d$ & $9.14$      & $9$       & $(3,0)$      & $1.5\%$  \\
$s$ & $182.78$    & $179$     & $(14,1)$     & $2.1\%$  \\
$c$ & $2495.11$   & $2491$    & $(29,14)$    & $0.16\%$ \\
$b$ & $8180.04$   & $8181$    & $(117,18)$   & $0.012\%$\\
$t$ & $338082$    & $338104$  & $(139,145)$  & $0.007\%$\\
\bottomrule
\end{tabular}
\end{table}

\section{Geometric Structure}
\label{sec:geometry}

\subsection{Spherical triangle of axis vectors}

The three axis vectors form a spherical triangle with side lengths
$48.1554^\circ$, $77.4835^\circ$, $68.4272^\circ$.
Interior angles (spherical law of cosines):
$A=72.13^\circ$, $B=92.36^\circ$, $C=49.68^\circ$,
giving spherical excess $E=34.18^\circ$.
The Cabibbo angle emerges from $\theta_{12}=48.16^\circ$ via the
Gaussian kernel construction, not from the excess.

\begin{figure}[htbp]
\centering
\includegraphics[width=\textwidth]{ckm_combined_figure.pdf}
\caption{Left: geodesic axis vectors on $S^2$ for the canonical triple
$\{\mathrm{aaB},\mathrm{AbA},\mathrm{AAb}\}$.
Right: spherical triangle with interior angles $72.1^\circ$,
$92.4^\circ$, $49.7^\circ$ and spherical excess $E=34.2^\circ$.}
\label{fig:geometry}
\end{figure}

\subsection{Arithmetic dichotomy with the PMNS sector}

The PMNS manifold $M_{\mathrm{PMNS}}=m003(-2,3)$ has imaginary trace
field $\mathbb{Q}(\sqrt{-3})$, giving loxodromic holonomy and large
CP violation ($\delta_{\mathrm{PMNS}}=195.91^\circ$, $0.55\%$ from
PDG~\cite{Gentry:Unified}). The real/imaginary trace field dichotomy
geometrically explains the CP asymmetry between quark and lepton sectors.

\section{Predictions and Open Problems}
\label{sec:predictions}

\paragraph{Falsifiable predictions.}
(1)~$|V_{us}|$ outside $[0.220,0.232]$ at $3\sigma$ falsifies the
construction.
(2)~$|J|>10^{-4}$ would require a twist mechanism absent from the
current framework.
(3)~A BSM quark mass ratio failing the $\mathbb{Z}[\sqrt{17}]$ norm
test to within $6\%$ breaks the arithmetic encoding.

\paragraph{Open problems.}
(1)~No geometric principle uniquely selects the word triple.
(2)~The $\mathbb{Z}[\sqrt{17}]$ norm encoding lacks a dynamical
derivation.
(3)~The $\mathrm{SU}(3)\times\mathrm{SU}(2)\times\mathrm{U}(1)$ gauge
structure is absent from the construction.

\section{Conclusions}
\label{sec:conclusions}

The compact hyperbolic 3-manifold $M_{\mathrm{CKM}}=m006(-5,2)$ encodes
the CKM matrix through two independent mechanisms:
\emph{arithmetic} (trace field $\mathbb{Q}(\sqrt{17})$, real, explaining
CP suppression via $[\mathrm{AbA}]=[\mathrm{AAb}]=2$) and
\emph{kinematic} (geodesic axis overlaps, Cabibbo angle error $0.19\%$,
improved fitness $0.00908$).
All six quark mass ratios approximate $\mathbb{Z}[\sqrt{17}]$ norms
with $p<0.002$. The framework is a geometric phenomenology with clearly
stated open problems and falsifiable predictions.

\section*{Data Availability}
\url{https://github.com/drmlgentry/hyperbolic-flavor-scan}

\section*{Acknowledgments}
The author thanks the developers of SnapPy and SageMath.

\begin{thebibliography}{99}
\bibitem{Cabibbo1963}
N.~Cabibbo, Phys.\ Rev.\ Lett.\ \textbf{10}, 531 (1963).
\bibitem{KobayashiMaskawa1973}
M.~Kobayashi and T.~Maskawa,
Prog.\ Theor.\ Phys.\ \textbf{49}, 652 (1973).
\bibitem{PDG2024}
S.~Navas et al.\ (PDG), Phys.\ Rev.\ D \textbf{110}, 030001 (2024).
\bibitem{SnapPy}
M.~Culler et al., \url{http://snappy.computop.org}.
\bibitem{Chinburg98}
T.~Chinburg et al.,
Ann.\ Scuola Norm.\ Sup.\ Pisa \textbf{30}, 1 (2001).
\bibitem{Gentry:Unified}
M.~L.~Gentry,
SSRN \href{https://doi.org/10.2139/ssrn.6775158}{6775158} (2026).
\end{thebibliography}

\end{document}